\documentclass[11pt]{article}
\usepackage[T1]{fontenc}
\usepackage[margin=1in]{geometry}
\usepackage{amsmath,amssymb,amsthm}
\usepackage{microtype}
\usepackage[hidelinks]{hyperref}
\newtheorem{theorem}{Theorem}
\newcommand{\R}{\mathbb{R}}
\begin{document}

\begin{center}
{\Large\bfseries Metrizability Cannot Be Removed from the Closed-Form Criterion}\par
\vspace{0.5em}
{\large A literature-implied negative answer to arXiv:1703.04883, Question 1}\par
\vspace{0.4em}
Run \texttt{fa\_banach\_001} --- lane 17 \qquad 13 August 2026
\end{center}

\medskip
\noindent\fbox{\begin{minipage}{0.94\linewidth}
\textbf{Status.} Full negative answer to the first loose end, obtained by applying a 1975 theorem of Dierolf. The metrizability hypothesis in Schmidt's Theorem~1.38 is essential: without it, assertions (i) and (ii) need not imply assertion (iii).
\end{minipage}}

\section{The source question}

Schmidt's \emph{Energy forms} \cite{Schmidt} proves the following criterion as Theorem~1.38. Let $(V,\tau)$ be a complete topological vector space and let $q$ be a quadratic form. If the form topology $(D(q),\tau_q)$ is metrizable, then any two of
\begin{enumerate}
\item[(i)] $q$ is closed;
\item[(ii)] $\ker q$ is $\tau$-closed and the canonical map
\[
(D(q)/\ker q,q)\longrightarrow (V/\ker q,\tau/\ker q)
\]
is continuous;
\item[(iii)] $(D(q)/\ker q,q)$ is a Hilbert space
\end{enumerate}
imply the third. Chapter~5, Section~5.1, Question~1 (printed page~163, arXiv PDF page~179) asks:

\begin{quote}\itshape
Does Theorem~1.38 hold without the assumption that $\tau_q$ is metrizable or are there counterexamples?
\end{quote}

There are counterexamples.

\section{The counterexample}

Let
\[
H_0=c_{00}=\{(x_n)_{n\geq1}:x_n=0\text{ for all but finitely many }n\}
\]
with the norm inherited from $\ell^2$. This is a Hausdorff locally convex space but is not complete.

Dierolf's Theorem~4 \cite[p.~76]{Dierolf} says that every topological vector space is a quotient of a complete Hausdorff topological vector space. Hence there are a complete Hausdorff topological vector space $(Z,\tau)$ and a quotient map
\[
\pi:Z\longrightarrow H_0.
\]
Write $L=\ker\pi$. Since $H_0\simeq Z/L$ is Hausdorff, $L$ is $\tau$-closed. Define the finite quadratic form
\[
q:Z\longrightarrow[0,\infty),\qquad q(z)=\|\pi z\|_2^2.
\]

\begin{theorem}
For this $q$, assertions \emph{(i)} and \emph{(ii)} of Schmidt's Theorem~1.38 hold, while assertion \emph{(iii)} fails.
\end{theorem}

\begin{proof}
We have $D(q)=Z$ and $\ker q=L$. The seminorm $q^{1/2}=\|\pi(\cdot)\|_2$ is $\tau$-continuous, because the quotient identification $Z/L\simeq H_0$ is a homeomorphism. The form topology is the smallest vector topology containing both $\tau$ and the topology induced by $q^{1/2}$; consequently
\[
\tau_q=\tau.
\]
Since $Z$ is complete, $(D(q),\tau_q)$ is complete. Thus $q$ is closed, proving (i).

We already observed that $L$ is closed. Under the quotient identification, both
\[
(D(q)/L,q)\quad\text{and}\quad(Z/L,\tau/L)
\]
are precisely $H_0$ with its $\ell^2$ norm topology. The canonical map in (ii) is therefore a topological isomorphism (indeed, an isometry on the form-norm side). This proves (ii).

Finally, $(D(q)/L,q)$ is the incomplete pre-Hilbert space $H_0$. For example, the finite truncations of $(1/n)_{n\geq1}\in\ell^2\setminus c_{00}$ form a Cauchy sequence in $H_0$ without a limit there. Hence (iii) fails.
\end{proof}

The example genuinely lies outside the omitted hypothesis. If $\tau_q=\tau$ were metrizable, then the quotient of the complete metrizable topological vector space $Z$ by its closed subspace $L$ would be complete, contradicting $Z/L\simeq H_0$.

\section{Provenance and scope}

This is a literature-implied answer: Dierolf's theorem predates the 2017 question by four decades but does not mention quadratic forms. The pullback construction above is the direct identification that closes the question. A bounded search through 13 August 2026 found the source question, Dierolf's theorem, and standard discussions of incomplete quotients of complete nonmetrizable spaces, but no explicit publication of this quadratic-form application. No priority claim is made.

The argument answers only Question~1 in Schmidt's loose-end chapter. It does not settle the later questions about Hilbert quotients of energy forms or local Poincar\'e inequalities.

\begin{sloppypar}
\begin{thebibliography}{9}
\raggedright
\bibitem{Schmidt}
M. Schmidt, \emph{Energy forms}, Ph.D. thesis, Friedrich Schiller University Jena, 2017; arXiv:1703.04883v1. See Theorem~1.38 and Chapter~5, Question~1.

\bibitem{Dierolf}
S. Dierolf, \emph{\"Uber Quotienten vollst\"andiger topologischer Vektorr\"aume}, Manuscripta Math. \textbf{17} (1975), 73--78. See Theorem~4, p.~76.
\end{thebibliography}
\end{sloppypar}

\end{document}
